\section{算法验证与性能分析}
\label{sec:ch5_validation}

\subsection{实验设置与评价指标}
\label{subsec:ch5_eval_setup}

感知与避障实验基于数字孪生闭环链路开展，实验数据由统一任务集与多轮重复试验汇总获得，图表采用一致统计口径生成。
评价指标覆盖识别性能、系统时延与避障决策质量三类。

识别任务采用 Precision/Recall/F1 的标准定义：
\begin{equation}
\mathrm{Precision}=\frac{TP}{TP+FP},\quad
\mathrm{Recall}=\frac{TP}{TP+FN},\quad
F_1=\frac{2PR}{P+R},
\label{eq:ch5_prf}
\end{equation}
并采用 AP 指标综合刻画精确率-召回率曲线下面积。
实时性指标采用单次前向耗时的均值、标准差与 $P95$。

\subsection{识别效果与方法对比}
\label{subsec:ch5_detection_compare}

SiaT-Hough 与主流分割方法的对比如表\ref{tab:ch5_method_compare}所示。

\begin{table}[H]
	\centering
	\small
	\caption{SiaT-Hough 与其他方法识别性能对比}
	\label{tab:ch5_method_compare}
	\begin{tabular}{p{0.28\textwidth}p{0.16\textwidth}p{0.16\textwidth}p{0.16\textwidth}}
		\toprule
		方法 & AP & mIoU & FPS \\
		\midrule
		SiaT-Hough & 0.81 & 0.77 & 31.2 \\
		Mask R-CNN & 0.74 & 0.69 & 12.8 \\
		YOLOv8-Seg & 0.76 & 0.71 & 44.5 \\
		\bottomrule
	\end{tabular}
\end{table}

可见，SiaT-Hough 在 AP 与 mIoU 上均优于对比方法，且在实时性上显著优于 Mask R-CNN；虽然 FPS 低于 YOLOv8-Seg，但其几何一致性更好，适合后续位姿估计和避障风险建模。

同时，在 CE4 局部视场重建实验中，系统以 $92\%$ 梯度阈值和 $28\%$ 阴影阈值获得 $3500$ 点岩石点云，并稳定提取 $8$ 个显著岩石候选（面积阈值 $20$ 像素），满足避障决策的目标级输入需求。

\subsection{系统实时性与避障收益分析}
\label{subsec:ch5_realtime_benefit}

针对系统级闭环效率，本文对关键模块执行重复基准测试，结果见表\ref{tab:ch5_runtime_benchmark}。

\begin{table}[H]
	\centering
	\small
	\caption{感知-规划关键模块运行时延统计}
	\label{tab:ch5_runtime_benchmark}
	\begin{tabular}{p{0.34\textwidth}p{0.18\textwidth}p{0.18\textwidth}p{0.18\textwidth}}
		\toprule
		模块 & 均值(ms) & 标准差(ms) & $P95$(ms) \\
		\midrule
		感知特征提取（30次） & 99.96 & 11.18 & 129.96 \\
		全局规划（20次） & 2487.57 & 72.60 & 2619.55 \\
		D3QN 单步前向（500次） & 1.21 & 0.39 & 1.96 \\
		\bottomrule
	\end{tabular}
\end{table}

进一步对全局路径平滑前后进行对比：路径长度由 $1289.19\,\mathrm{m}$ 降至 $1280.99\,\mathrm{m}$（下降 $0.64\%$），累计转向量由 $79.33$ 降至 $50.96$（下降 $35.76\%$）。
这说明在不显著增加计算开销的前提下，风险引导的路径平滑可有效降低机动剧烈度，提高执行稳定性。

在局部避障策略层面，三类候选策略风险评分分别为左绕 $0.741215$、右绕 $0.785823$、等待 $0.789683$，系统最终选择左绕路径，验证了式\eqref{eq:ch5_policy_select}的有效性。

\subsection{消融实验与模块贡献度分析}
\label{subsec:ch5_ablation}

为进一步验证增强型 ORB-SLAM2 中各子模块的必要性，本节在统一数据划分与统一评价口径下开展消融实验。设完整模型记为
\begin{equation}
\mathcal{F}_{full}=\mathcal{F}_{SiaT}\circ\mathcal{F}_{Hough}\circ\mathcal{F}_{temp}\circ\mathcal{F}_{geo},
\label{eq:ch5_full_pipeline}
\end{equation}
其中 $\mathcal{F}_{temp}$ 表示时序一致性约束，$\mathcal{F}_{geo}$ 表示几何一致性筛选。对任一移除模块 $i$，其相对贡献度定义为
\begin{equation}
\eta_i=\frac{\mathcal{M}(\mathcal{F}_{full})-\mathcal{M}(\mathcal{F}_{-i})}{\mathcal{M}(\mathcal{F}_{full})}\times 100\%,
\label{eq:ch5_contribution}
\end{equation}
其中 $\mathcal{M}(\cdot)$ 表示目标评价指标（本文以 $F_1$ 为主指标）。

\begin{table}[H]
	\centering
	\small
	\caption{SiaT-Hough 关键模块消融结果}
	\label{tab:ch5_ablation}
	\begin{tabular}{p{0.36\textwidth}p{0.12\textwidth}p{0.12\textwidth}p{0.12\textwidth}p{0.16\textwidth}}
		\toprule
		配置 & AP & Recall & $F_1$ & 计数 MAE \\
		\midrule
		完整模型（SiaT+Hough+时序+几何） & 0.81 & 0.84 & 0.82 & 0.73 \\
		去除 SiaT 语义增强 & 0.75 & 0.79 & 0.76 & 1.21 \\
		去除 Hough 几何投票 & 0.74 & 0.77 & 0.75 & 1.35 \\
		去除时序一致性约束 & 0.78 & 0.80 & 0.79 & 0.98 \\
		去除几何一致性筛选 & 0.77 & 0.82 & 0.79 & 1.04 \\
		\bottomrule
	\end{tabular}
\end{table}

表\ref{tab:ch5_ablation}表明：SiaT 与 Hough 子模块对识别与计数精度均具有显著贡献，其中 Hough 对降低伪匹配点导致的计数偏差尤为关键；时序与几何约束则主要提升结果稳定性，避免单帧异常误检对局部避障决策造成放大影响。

\begin{algorithm}[H]
	\caption{SiaT-Hough 前端融合伪代码}
	\label{alg:ch5_siathough_fusion}
	\begin{algorithmic}[1]
		\REQUIRE 输入帧 $I_t$，历史关键帧集合 $\mathcal{K}_{t-1}$
		\ENSURE 岩石候选集 $\mathcal{R}_t$
		\STATE 计算语义增强特征 $\mathbf{f}_t\leftarrow \mathcal{F}_{SiaT}(I_t)$
		\STATE 生成几何投票集合 $\mathcal{V}_t\leftarrow \mathcal{F}_{Hough}(\mathbf{f}_t)$
		\STATE 执行时序一致性筛选 $\tilde{\mathcal{V}}_t\leftarrow \mathcal{F}_{temp}(\mathcal{V}_t,\mathcal{K}_{t-1})$
		\STATE 执行几何一致性筛选 $\mathcal{R}_t\leftarrow \mathcal{F}_{geo}(\tilde{\mathcal{V}}_t)$
		\STATE 将 $\mathcal{R}_t$ 回写至地图并更新关键帧缓存
	\end{algorithmic}
\end{algorithm}

\begin{figure}[H]
	\centering
	\resizebox{0.72\textwidth}{!}{%
	\begin{tikzpicture}[font=\small, >=Stealth]
		\node[draw, rounded corners, fill=blue!10, align=center, minimum width=6.8cm, minimum height=1.0cm] (a) at (5,8.8) {输入图像与关键帧缓存};
		\node[draw, rounded corners, fill=cyan!12, align=center, minimum width=6.8cm, minimum height=1.0cm] (b) at (5,7.0) {SiaT 语义增强特征提取};
		\node[draw, rounded corners, fill=green!12, align=center, minimum width=6.8cm, minimum height=1.0cm] (c) at (5,5.2) {Hough 几何投票与候选聚类};
		\node[draw, rounded corners, fill=orange!14, align=center, minimum width=6.8cm, minimum height=1.0cm] (d) at (5,3.4) {时序一致性与几何一致性筛选};
		\node[draw, rounded corners, fill=purple!10, align=center, minimum width=6.8cm, minimum height=1.0cm] (e) at (5,1.6) {输出岩石候选并回写地图};
		\draw[->, thick, draw=black!75] (a) -- (b);
		\draw[->, thick, draw=black!75] (b) -- (c);
		\draw[->, thick, draw=black!75] (c) -- (d);
		\draw[->, thick, draw=black!75] (d) -- (e);
		\draw[->, thick, draw=blue!70!black] (e.west) .. controls (1.2,1.1) and (1.2,7.8) .. node[left, align=center]{阈值自适应\\参数更新} (b.west);
	\end{tikzpicture}%
	}
	\caption{SiaT-Hough 前端融合程序框图}
	\label{fig:ch5_siathough_flow}
\end{figure}

图\ref{fig:ch5_siathough_flow}给出了前端融合链路的核心步骤与反馈关系：语义增强负责提升岩石候选的可分性，几何投票负责抑制局部纹理伪匹配，筛选与回写环节保证识别结果可持续用于后续重建与避障决策。

\subsection{误差来源分解与不确定性传播分析}
\label{subsec:ch5_uncertainty}

考虑识别-重建-计数链路的误差耦合，本节将总误差表示为
\begin{equation}
\varepsilon_{total}=\varepsilon_{det}+\varepsilon_{match}+\varepsilon_{tri}+\varepsilon_{count},
\label{eq:ch5_error_decompose}
\end{equation}
其中 $\varepsilon_{det}$、$\varepsilon_{match}$、$\varepsilon_{tri}$、$\varepsilon_{count}$ 分别对应检测、匹配、三角化和计数阶段误差。对小扰动条件下的线性化传播，设状态向量为 $\mathbf{x}$，观测向量为 $\mathbf{z}$，则协方差传播近似为
\begin{equation}
\mathbf{\Sigma}_{k+1}\approx \mathbf{J}_f\mathbf{\Sigma}_k\mathbf{J}_f^\top+\mathbf{Q},
\label{eq:ch5_cov_prop}
\end{equation}
其中 $\mathbf{J}_f$ 为状态转移雅可比矩阵，$\mathbf{Q}$ 为过程噪声协方差。

\begin{table}[H]
	\centering
	\small
	\caption{计数误差来源分解与相对贡献}
	\label{tab:ch5_error_source}
	\begin{tabular}{p{0.34\textwidth}p{0.18\textwidth}p{0.18\textwidth}p{0.18\textwidth}}
		\toprule
		误差来源 & 均值偏差 & 方差贡献率 & 工程影响等级 \\
		\midrule
		弱纹理引发的漏检 & 0.31 & 28.4\% & 高 \\
		阴影边界误检 & 0.24 & 22.1\% & 中高 \\
		跨视角匹配漂移 & 0.18 & 19.7\% & 中 \\
		三角化深度抖动 & 0.15 & 16.3\% & 中 \\
		后处理阈值波动 & 0.09 & 13.5\% & 中低 \\
		\bottomrule
	\end{tabular}
\end{table}

由表\ref{tab:ch5_error_source}可见，识别阶段（漏检与误检）贡献了主要误差方差，说明继续提高前端语义与几何联合判别能力仍是提升计数精度的首要方向。三角化与后处理误差占比相对较低，表明当前重建后端具备较好稳定性。

\begin{figure}[H]
	\centering
	\resizebox{0.70\textwidth}{!}{%
	\begin{tikzpicture}[font=\small, >=Stealth]
		\node[draw, rounded corners, fill=red!10, align=center, minimum width=3.6cm, minimum height=1.0cm] (d1) at (3.0,8.2) {检测误差\\$\varepsilon_{det}$};
		\node[draw, rounded corners, fill=orange!14, align=center, minimum width=3.6cm, minimum height=1.0cm] (d2) at (7.0,8.2) {匹配误差\\$\varepsilon_{match}$};
		\node[draw, rounded corners, fill=yellow!20, align=center, minimum width=4.2cm, minimum height=1.0cm] (d3) at (5.0,5.8) {三角化误差\\$\varepsilon_{tri}$};
		\node[draw, rounded corners, fill=blue!10, align=center, minimum width=4.2cm, minimum height=1.0cm] (d4) at (5.0,3.5) {计数误差\\$\varepsilon_{count}$};
		\node[draw, rounded corners, fill=purple!10, align=center, minimum width=4.6cm, minimum height=1.0cm] (d5) at (5.0,1.4) {风险映射与避障代价修正};
		\draw[->, thick, draw=black!75] (d1) -- (d3);
		\draw[->, thick, draw=black!75] (d2) -- (d3);
		\draw[->, thick, draw=black!75] (d3) -- (d4);
		\draw[->, thick, draw=black!75] (d4) -- (d5);
		\draw[->, thick, draw=blue!70!black] (d5.west) .. controls (1.4,1.2) and (1.4,7.4) .. node[left, align=center]{不确定性\\回传} (d1.west);
	\end{tikzpicture}%
	}
	\caption{识别-重建-计数误差传播架构图}
	\label{fig:ch5_error_chain}
\end{figure}

针对图\ref{fig:ch5_error_chain}所示误差链路，可将计数误差上界写为
\begin{equation}
|\varepsilon_{count}|\le
\lambda_1|\varepsilon_{det}|+
\lambda_2|\varepsilon_{match}|+
\lambda_3|\varepsilon_{tri}|,
\label{eq:ch5_count_error_bound}
\end{equation}
其中 $\lambda_1,\lambda_2,\lambda_3$ 为链路灵敏度系数。实验统计显示 $\lambda_1>\lambda_2>\lambda_3$，与前述方差贡献结论一致。

\subsection{跨场景泛化能力与鲁棒性评估}
\label{subsec:ch5_generalization}

为验证方法在不同月面工况下的稳定性，本节从光照条件、地形复杂度与岩石密度三个维度构建对比实验。设场景域记为 $\mathcal{D}_s$，模型在目标域 $\mathcal{D}_t$ 上的性能保持率定义为
\begin{equation}
\rho_{gen}=\frac{\mathcal{M}_{\mathcal{D}_t}}{\mathcal{M}_{\mathcal{D}_s}}.
\label{eq:ch5_generalization_ratio}
\end{equation}

\begin{table}[H]
	\centering
	\small
	\caption{跨光照与跨地形条件下的识别鲁棒性对比}
	\label{tab:ch5_cross_scene}
	\begin{tabular}{p{0.28\textwidth}p{0.14\textwidth}p{0.14\textwidth}p{0.14\textwidth}p{0.14\textwidth}}
		\toprule
		场景条件 & Precision & Recall & $F_1$ & $\rho_{gen}$ \\
		\midrule
		基准光照+中等地形 & 0.83 & 0.84 & 0.83 & 1.000 \\
		低照度+阴影遮挡增强 & 0.79 & 0.81 & 0.80 & 0.964 \\
		高反差直射光 & 0.78 & 0.80 & 0.79 & 0.952 \\
		起伏地形+岩石高密度 & 0.77 & 0.79 & 0.78 & 0.940 \\
		低照度+起伏地形耦合 & 0.75 & 0.77 & 0.76 & 0.916 \\
		\bottomrule
	\end{tabular}
\end{table}

结果显示，算法在多种工况下均保持较高识别稳定性，性能保持率最低为 $0.916$。在耦合极端场景中，性能下降主要由阴影边界与高起伏地形共同导致的局部纹理退化引起，但整体仍满足避障任务对岩石先验输入的稳定性要求。

\begin{figure}[H]
	\centering
	\resizebox{0.70\textwidth}{!}{%
	\begin{tikzpicture}[font=\small, >=Stealth]
		\node[draw, rounded corners, fill=blue!10, align=center, minimum width=3.8cm, minimum height=1.0cm] (s1) at (5,8.5) {场景组构建\\（光照/地形/密度）};
		\node[draw, rounded corners, fill=cyan!12, align=center, minimum width=3.8cm, minimum height=1.0cm] (s2) at (5,6.5) {统一推理与指标计算};
		\node[draw, rounded corners, fill=green!12, align=center, minimum width=3.8cm, minimum height=1.0cm] (s3) at (5,4.5) {统计显著性检验};
		\node[draw, rounded corners, fill=orange!14, align=center, minimum width=3.8cm, minimum height=1.0cm] (s4) at (5,2.5) {误差回溯与主因定位};
		\node[draw, rounded corners, fill=purple!10, align=center, minimum width=3.8cm, minimum height=1.0cm] (s5) at (5,0.9) {阈值修订与再评估};
		\draw[->, thick, draw=black!75] (s1) -- (s2);
		\draw[->, thick, draw=black!75] (s2) -- (s3);
		\draw[->, thick, draw=black!75] (s3) -- (s4);
		\draw[->, thick, draw=black!75] (s4) -- (s5);
		\draw[->, thick, draw=blue!70!black] (s5.east) .. controls (9.3,1.0) and (9.3,7.7) .. node[right, align=center]{闭环迭代} (s1.east);
	\end{tikzpicture}%
	}
	\caption{跨场景鲁棒性验证流程图}
	\label{fig:ch5_cross_scene_flow}
\end{figure}

图\ref{fig:ch5_cross_scene_flow}对应的流程强调“场景构建-统一评估-显著性检验-误差回溯-参数修订”的闭环机制，可有效降低由单一实验工况造成的性能高估风险，提升结论的统计可信度。

\subsection{工程部署评估与可复现性分析}
\label{subsec:ch5_engineering}

围绕工程落地需求，本节补充算法复杂度、资源消耗和复现实验稳定性分析。完整管线单帧计算复杂度可近似写为
\begin{equation}
\mathcal{O}_{frame}\approx \mathcal{O}(N_f\log N_f)+\mathcal{O}(N_v)+\mathcal{O}(N_p),
\label{eq:ch5_complexity}
\end{equation}
其中 $N_f$ 为候选特征点数，$N_v$ 为几何投票数，$N_p$ 为后处理点云规模。

\begin{table}[H]
	\centering
	\small
	\caption{感知链路部署开销统计}
	\label{tab:ch5_deployment}
	\begin{tabular}{p{0.30\textwidth}p{0.18\textwidth}p{0.18\textwidth}p{0.20\textwidth}}
		\toprule
		指标 & 平均值 & 峰值 & 备注 \\
		\midrule
		前端推理时延（ms） & 99.96 & 136.42 & 满足在线处理 \\
		点云重建时延（ms） & 42.73 & 65.11 & 与候选数近线性相关 \\
		内存占用（MB） & 412.8 & 569.4 & 峰值出现在重建阶段 \\
		单任务能耗（归一化） & 1.00 & 1.17 & 与地形复杂度相关 \\
		\bottomrule
	\end{tabular}
\end{table}

进一步地，本文在相同实验协议下进行多轮重复试验，关键指标（$F_1$、计数 MAE、端到端时延）变异系数均低于 $5\%$，表明系统在工程实现层面具备良好的稳定性与可复现性。该结论支撑后续章节在规划任务中继续调用本章感知输出，形成可靠的跨章节技术闭环。

\subsection{极端工况案例分析与失效模式讨论}
\label{subsec:ch5_failure_modes}

为避免仅在平均指标层面给出结论，本节补充极端工况下的案例分析。考虑低照度、高反差、密集岩石、强起伏四类高风险场景，定义单次任务失效率为
\begin{equation}
P_{fail}=1-\frac{N_{safe}}{N_{total}},
\label{eq:ch5_fail_prob}
\end{equation}
其中 $N_{safe}$ 表示满足“无碰撞、无越界、关键目标未漏检”三条件的任务数。进一步定义风险敏感增益
\begin{equation}
G_{risk}=\frac{R_{base}-R_{enh}}{R_{base}},
\label{eq:ch5_risk_gain}
\end{equation}
其中 $R_{base}$ 与 $R_{enh}$ 分别为基线模型与增强模型的综合风险评分。$G_{risk}>0$ 表示增强策略有效降低风险。

\begin{table}[H]
	\centering
	\small
	\caption{极端工况下的识别与避障联合表现}
	\label{tab:ch5_extreme_cases}
	\begin{tabular}{p{0.30\textwidth}p{0.14\textwidth}p{0.14\textwidth}p{0.14\textwidth}p{0.16\textwidth}}
		\toprule
		工况类型 & $F_1$ & 计数 MAE & $P_{fail}$ & $G_{risk}$ \\
		\midrule
		低照度+长阴影 & 0.74 & 1.42 & 0.117 & 0.236 \\
		高反差直射光 & 0.76 & 1.36 & 0.103 & 0.221 \\
		岩石高密度聚集 & 0.72 & 1.58 & 0.132 & 0.248 \\
		起伏地形+稀疏纹理 & 0.73 & 1.47 & 0.124 & 0.231 \\
		多因素耦合极端场景 & 0.69 & 1.83 & 0.158 & 0.264 \\
		\bottomrule
	\end{tabular}
\end{table}

由表\ref{tab:ch5_extreme_cases}可知，在多因素耦合场景中系统性能下降最为明显，但增强策略仍能保持正风险增益，说明该方法对极端工况具有一定容错性。造成性能下滑的主因是“纹理弱化+投影畸变+遮挡叠加”共同触发的候选不稳定，且该问题具有明显帧间传播特性。

\begin{figure}[H]
	\centering
	\resizebox{0.72\textwidth}{!}{%
	\begin{tikzpicture}[font=\small, >=Stealth]
		\node[draw, rounded corners, fill=red!12, align=center, minimum width=7.0cm, minimum height=1.0cm] (f1) at (5,8.6) {极端场景触发与风险预警};
		\node[draw, rounded corners, fill=orange!14, align=center, minimum width=7.0cm, minimum height=1.0cm] (f2) at (5,6.8) {异常帧识别与候选稳定性检测};
		\node[draw, rounded corners, fill=yellow!20, align=center, minimum width=7.0cm, minimum height=1.0cm] (f3) at (5,5.0) {误差分解与主因归类};
		\node[draw, rounded corners, fill=green!12, align=center, minimum width=7.0cm, minimum height=1.0cm] (f4) at (5,3.2) {阈值重整与策略切换};
		\node[draw, rounded corners, fill=blue!10, align=center, minimum width=7.0cm, minimum height=1.0cm] (f5) at (5,1.4) {闭环复测与稳定性确认};
		\draw[->, thick, draw=black!75] (f1) -- (f2);
		\draw[->, thick, draw=black!75] (f2) -- (f3);
		\draw[->, thick, draw=black!75] (f3) -- (f4);
		\draw[->, thick, draw=black!75] (f4) -- (f5);
		\draw[->, thick, draw=blue!70!black] (f5.east) .. controls (9.2,1.2) and (9.2,7.8) .. node[right, align=center]{在线\\迭代} (f1.east);
	\end{tikzpicture}%
	}
	\caption{极端工况失效诊断与修正程序框图}
	\label{fig:ch5_failure_diagnosis_flow}
\end{figure}

图\ref{fig:ch5_failure_diagnosis_flow}展示了极端工况下的诊断闭环。该闭环与常规评估流程相比，增加了“异常帧稳定性检测”和“主因归类”两个关键环节，使系统能够在不改变总体架构的前提下进行局部参数重整，从而抑制错误在感知-规划链路中的级联放大。

在工程实践中，可将失效模式划分为三类：第一类为检测层失效（漏检、误检）；第二类为几何层失效（匹配漂移、深度偏置）；第三类为决策层失效（安全边界估计偏差）。针对三类失效分别配置“阈值回退、候选重评分、安全距离保守扩张”机制，可将高风险任务的失败率进一步控制在可接受范围内。该结论对月面长周期任务尤为重要，因为长周期任务对单次局部失败的累积敏感性显著高于短程验证任务。

\section{本章小结}
\label{sec:ch5_summary}

本章围绕“数字孪生感知如何服务可执行避障”展开，形成了从识别到决策的完整技术链路。
核心贡献可概括为三点：
\begin{myitemize}
	\item 提出并实现了增强型 ORB-SLAM2 前端，通过 SiaT-Hough 模块将语义一致性与霍夫几何投票融合，提升了月面复杂光照与重复纹理条件下的匹配鲁棒性；
	\item 构建了岩石特征提取、三维重建与目标计数流程，完成了从二维检测结果到三维障碍先验的转换，为后续规划提供结构化输入；
	\item 设计了风险加权避障策略与动态安全距离模型，将沉陷、碰撞、能耗与时延约束统一到同一决策代价函数中，并通过实验验证了策略有效性。
\end{myitemize}

研究局限主要在于：当前计数误差仍依赖有限人工核验，且真实月面标注数据规模不足；此外，端到端系统在规划环节仍存在秒级耗时瓶颈。
后续研究将扩展高质量标注集并引入轻量化并行推理机制，进一步提升感知-决策闭环效率与工程可部署性。
